\chapter*{Kết luận}
\addcontentsline{toc}{chapter}{Kết luận}

Đề tài đã tập trung giải quyết bài toán xây dựng một ứng dụng di động hỗ trợ tập luyện và theo dõi sức khỏe cá nhân hóa trên nền tảng Android. Trên cơ sở phân tích nhu cầu thực tế của người dùng phổ thông, hệ thống Fitty được thiết kế để không chỉ cung cấp một danh sách bài tập đơn lẻ mà còn tạo ra một luồng trải nghiệm hoàn chỉnh gồm xác thực truy cập, onboarding, sinh kế hoạch tập khởi đầu, tra cứu bài tập, thực hiện buổi tập, theo dõi tiến độ và duy trì tương tác qua thông báo.

Về mặt kỹ thuật, hệ thống đã hiện thực được một kiến trúc nhiều tầng tương đối rõ ràng với giao diện Jetpack Compose, điều hướng bằng Navigation Compose, dependency injection bằng Hilt, lưu trữ cục bộ bằng Room và DataStore, cùng các dịch vụ phía sau dựa trên Firebase. Một trong những điểm nổi bật của hệ thống là việc áp dụng mô hình offline-first cho phân hệ bài tập, qua đó tách metadata và media, đồng bộ dữ liệu qua repository chuyên trách và sử dụng Room như nguồn dữ liệu ổn định cho giao diện. Bên cạnh đó, quy trình onboarding gắn với logic sinh starter plan và use case hoàn thành buổi tập gắn với cập nhật summary, trạng thái kế hoạch và thống kê người dùng là các điểm thể hiện rõ tính liên kết giữa các phân hệ của hệ thống.

Qua quá trình khảo sát mã nguồn, tài liệu kỹ thuật và các kế hoạch triển khai hiện có, có thể đánh giá rằng hệ thống đã đáp ứng phần lớn những chức năng quan trọng của một ứng dụng hỗ trợ tập luyện cá nhân hóa. Các phân hệ cốt lõi như xác thực, onboarding, kế hoạch tập, nội dung bài tập, workout session, tracking, profile, thông báo cục bộ và FCM đều đã được hiện thực ở mức có thể vận hành như một nguyên mẫu chức năng có cấu trúc rõ ràng. Ngoài ra, việc có sẵn các unit test cho nhiều use case và repository quan trọng cũng là cơ sở tốt để tiếp tục mở rộng và nâng cao chất lượng hệ thống.

Tuy nhiên, hệ thống vẫn còn một số điểm cần tiếp tục hoàn thiện trong tương lai. Một số chỉ số theo dõi hiện vẫn được ước lượng từ dữ liệu sẵn có thay vì được hỗ trợ hoàn toàn từ một mô hình nghiệp vụ sâu hơn. Chất lượng trải nghiệm ở các màn hình tracking và buổi tập nhanh còn có thể nâng cấp thêm nếu dữ liệu động và các cấu hình persisted được làm giàu hơn. Bên cạnh đó, việc tăng cường kiểm thử tích hợp trên thiết bị thật và tiếp tục chuẩn hóa nội dung quản lý từ Firebase sẽ giúp hệ thống tiến gần hơn tới một sản phẩm có mức hoàn thiện cao.

Trong các hướng phát triển tiếp theo, hệ thống có thể được mở rộng theo nhiều hướng như cá nhân hóa kế hoạch tập sâu hơn, bổ sung biểu đồ theo dõi trực quan, nâng cấp phân tích dinh dưỡng và thể trạng, tăng cường đồng bộ hóa nội dung hướng dẫn, hoặc tích hợp thêm các nguồn dữ liệu sức khỏe. Với kiến trúc hiện tại, Fitty đã có nền tảng phù hợp để tiếp tục được phát triển thành một hệ thống hoàn chỉnh hơn trong phạm vi sản phẩm thực tế.

Tóm lại, đề tài đã đạt được mục tiêu xây dựng một hệ thống ứng dụng di động hỗ trợ tập luyện và theo dõi sức khỏe cá nhân hóa với kiến trúc hợp lý, phạm vi chức năng rõ ràng và khả năng mở rộng tốt. Các kết quả đạt được từ phân tích, thiết kế và cài đặt hệ thống là cơ sở có giá trị cho các bước phát triển tiếp theo của ứng dụng cũng như cho việc nghiên cứu và hiện thực các bài toán tương tự trong lĩnh vực ứng dụng sức khỏe trên thiết bị di động.
